% Del encendido a la partida: la lista numerada de pantallas, en orden.
\begin{enumerate}
\item \textbf{Mensaje en la pantalla del MSX.} Durante unos dos segundos la pantalla propia del MSX muestra
\texttt{CAPCOM 1985}, \texttt{DIHALT STUDIO 2026}, \texttt{VERSION \portVersion} y \texttt{V9990 found!}. Si pone \texttt{THIS GAME RUNS ONLY ON MSX TURBO
R} (el ordenador no es un MSX turbo R) o \texttt{V9990 NOT found!!!} (no hay V9990), el juego se detiene ahí.
\item \textbf{Logo de DiHalt.} Con MoonSound o Wozblaster, el logo permanece en el V9990 mientras suena la fanfarria de DiHalt Studio, de diez segundos. Con otras configuraciones se muestra durante cuatro segundos en silencio.
\item \textbf{Portada.} La ilustración a pantalla completa. Espera hasta que pulses \key{SPACE} o el botón de disparo del
joystick.
\item \textbf{Mensaje de homenaje.} Se escribe en pantalla un breve texto: este juego es un homenaje al Commando original
de Capcom.
\item \textbf{Título.} La pantalla con las puntuaciones de bonificación y \texttt{PRESS FIRE TO START} parpadeando, y
después el gran logo de \textsc{Commando} con el ranking de las siete mejores puntuaciones. Se repite hasta que empieces.
\item \textbf{Empezar.} Pulsa \key{1} o disparo para una partida normal, o \key{2} o el botón de granada para una partida
fácil (las balas enemigas van más despacio). No hay monedas ni créditos, y el juego es para un jugador.
\item \textbf{La partida.} Un helicóptero te deja al principio del área~1 y empieza el juego.
\end{enumerate}

Cuando pierdes la última vida aparece \texttt{GAME OVER}. Si tu puntuación está entre las siete mejores puedes escribir tu
nombre (sección~8); después se muestra el ranking y vuelven las pantallas del título.
